\section{Introduction}
\label{sec:intro}

Large language models (LLMs) now support question answering over
enterprise documents, regulatory and financial assistants, and
clinical decision support
\citep{lewis2020rag,gao2023alce,yang2024cragbench,aarab2025bank,
yang2025integrating}. Because parametric
knowledge is incomplete, can become stale, and may produce
hallucinations, retrieval-augmented generation (RAG) grounds model
outputs in external sources: relational tables, text corpora under
sparse and dense indexes, vector stores, and structured knowledge APIs
\citep{lewis2020rag,chen2020hybridqa}. A production RAG service is a
decision system, not merely a retriever followed by a generator. It may
combine several knowledge sources, rerankers, verifiers, and LLMs whose
prices differ by two orders of magnitude. For each request, an operator
must decide how much evidence to acquire and when an inexpensive answer
is sufficient. That decision must respect limits on wrong,
unsupported, stale, or late answers. Average development-set accuracy
does not provide such an operational guarantee.

\begin{figure*}[t]
\centering
\includegraphics[width=0.88\textwidth]{figures/fig_paradigm.png}
\caption{Three deployment paradigms for RAG knowledge services.
(a)~Point-estimate tuning deploys the cheapest configuration whose
\emph{measured} risk clears the target --- an uncorrected
multiple-testing procedure with no feasibility guarantee.
(b)~Distribution-free risk control bounds the risk of a \emph{fixed}
pipeline but neither searches a configuration space nor reduces cost,
and is silent under drift. (c)~\system\ combines search,
certification, and monitoring: an evidence contract is certified by a fixed-sequence
test over a policy space ordered by risk vouchers, and an anytime-valid
e-process monitor closes the loop after deployment.}
\label{fig:paradigm}
\end{figure*}

Current deployment practice does not support this decision.
Practitioners sample a development set, estimate each configuration's
quality, and choose the cheapest configuration whose \emph{estimate}
clears the target, an approach we call point-estimate tuning
(Figure~\ref{fig:paradigm}a), shared
by workload-level pipeline optimizers \citep{russo2025abacus,
liu2025palimpzest,patel2024lotus} and by adaptive RAG routers and LLM
cascades \citep{jeong2024adaptiverag,chen2023frugalgpt,ong2024routellm}.
This selection is an uncorrected multiple-testing procedure. It favors
configurations whose measurement noise happens to be favorable, and the
problem worsens as the candidate set grows. As we quantify in
Section~\ref{sec:results}, state-of-practice optimizers deploy policies
whose true violation rate exceeds the declared target in up to 64\% of
calibration draws. Even a correctly tuned pipeline offers no protection after deployment,
when workload mix shifts, indexes fall behind the world, and knowledge
APIs degrade --- the knowledge drift problem that static evaluation
cannot see.

Distribution-free risk control addresses the statistical problem
(Figure~\ref{fig:paradigm}b): conformal
and Learn-then-Test methods \citep{angelopoulos2021ltt,
angelopoulos2024crc} can bound the violation rate of a \emph{fixed}
pipeline with finite-sample confidence, and recent work has applied
them to RAG abstention and answer sets \citep{li2024traq,kang2024crag}.
These methods treat the pipeline as given: they neither search a
configuration space nor reduce cost, and they are silent about what
happens when the calibrated distribution drifts. Conversely, the
systems that do search --- routers, cascades, Bayesian optimization ---
provide no feasibility statement. Search, certification, and
post-deployment monitoring therefore need to be designed together so
that cost optimization does not weaken the statistical guarantee.

We propose \system, a risk-aware expert system for RAG operations
(Figure~\ref{fig:paradigm}c). A domain expert states the operating
requirements as an \emph{evidence contract}: admissible violation rates
for answer quality, citation completeness, evidence freshness, and
latency, together with a confidence budget $\delta$. \system translates
these requirements into an auditable recommendation. It either returns
a certified policy or reports that the current system cannot meet the
contract. It never hides infeasibility behind a favorable point
estimate.

The system has three decision stages. First, it compiles a lattice of
knowledge-access plans using retrieval truncation, reranker removal,
pre-filter pushdown, and generator downgrade. Empirical
\emph{risk vouchers} summarize how each rewrite has affected prior
answers. The vouchers compose without independence assumptions and
guide the search, but they do not authorize deployment. Second,
\system turns the plans into progressive policies. Each policy starts
cheaply and acquires stronger evidence while a learned sufficiency
score remains low. The candidate family also includes utility routers,
group-conditional variants, and human deferral. A fixed-sequence
Hoeffding--Bentkus test certifies the least-cost policy in the accepted
prefix. Third, an anytime-valid e-detector audits the running service.
It detects drift, identifies the data that should be used for
recalibration, and restarts the decision process. A geometric error
schedule keeps this loop valid over repeated alarms.

This work provides an end-to-end decision-support architecture for RAG
over heterogeneous knowledge sources. Its main contributions are:
\begin{itemize}
\item We define evidence contracts as an operator-facing interface for
quality, citation, freshness, latency, and confidence requirements
(Section~\ref{sec:prelim}).
\item We develop a policy-recommendation engine that combines a
rewrite-generated plan lattice, composable vouchers, learned
sufficiency scores, and fixed-sequence certification. It recommends a
least-cost certified policy or declines deployment
(Section~\ref{sec:method}).
\item We close the operating loop with an anytime-valid drift monitor,
sampled auditing, changepoint localization, safe-mode escalation, and
repeated recertification under a lifetime error budget
(Section~\ref{sec:monitor}).
\item We evaluate the full system on four public benchmarks, four LLM
families, and two serving backends. The study covers reliability,
cost, strict and joint contracts, infeasibility, audit rates, and
drift response (Sections~\ref{sec:setup}--\ref{sec:results}).
\end{itemize}

The remainder of this paper is organized as follows.
Section~\ref{sec:related} reviews related work.
Section~\ref{sec:prelim} formalizes evidence contracts and the
certified deployment problem. Section~\ref{sec:method} presents the
\system\ framework. Sections~\ref{sec:setup} and~\ref{sec:results}
report the experimental setup and results. Section~\ref{sec:discussion}
discusses implications, limitations, and threats to validity, and
Section~\ref{sec:conclusion} concludes.
